\subsection{Topology-preserving loss}
\label{sec:met:loss}

\paragraph{Motivation.}
Two properties of neuronal morphology make our recordings difficult to
segment: branches can be only a few pixels wide, and neighbouring branches
often touch or overlap due to the image resolution
\cite{grim2025supervoxel}. An error at this scale --- negligible under a
per-pixel or region-overlap loss --- can therefore split one branch into
two, or merge two branches into one. The loss score is almost unchanged, yet
the reconstructed morphology is no longer usable for the measurements that
motivate the segmentation.

Splits and merges are topological errors: they change the connectivity of
the segmented neuron while leaving its area essentially intact, and neither
a per-pixel nor a region-overlap term can register them. We therefore
include two terms that operate on topology. cbDice \cite{shi2024cbdice}
compares the centrelines of prediction and target rather than their areas,
and weights that comparison by local thickness so that a thin branch counts
more than the cell body --- exactly where an area-based loss is weakest. The
supervoxel loss \cite{grim2025supervoxel} examines the connected components
of the error mask and identifies those whose addition or removal would
change the number of connected components in the prediction; it penalises
false splits and false merges directly, which cbDice, comparing centrelines,
does not.

\paragraph{The composite loss.}
% NOTE (verified from gpu-checkpoints/training_log.json): loss_iou and
% loss_class are identically 0.0 across every logged epoch - they are gated
% on pred_obj_scores, which the training wrapper does not emit. The loss
% actually optimised has FOUR terms. Do not reinstate the six-term form.
We supervise the predicted mask with a weighted sum of four terms --- two
that measure agreement in the conventional way, per pixel and per region,
and two that measure topology:
\begin{equation}
\mathcal{L} = 5\,\mathcal{L}_{\text{focal}} + 2\,\mathcal{L}_{\text{Dice}}
 + 5\,\mathcal{L}_{\text{cbDice}} + 0.3\,\mathcal{L}_{\text{sv}} .
\label{eq:loss}
\end{equation}
The first pair is the mask supervision under which SAM and SAM\,2 were
themselves trained \cite{kirillov2023sam,ravi2025sam2}; we keep it so that
adaptation proceeds under the same supervision that produced the pretrained
weights, and retune only the relative weights. Focal loss
\cite{lin2017focal} is a cross-entropy in which every pixel is scaled by a
class factor $\alpha_t$ and by $(1-p_t)^{\gamma}$, so that a pixel already
classified confidently contributes almost nothing; we use the standard
pairing $\gamma = 2.0$, $\alpha = 0.25$, the focusing term having already
absorbed most of the imbalance that $\alpha$ would otherwise correct. The
target is a small share of each frame and that share varies by an order of
magnitude across recordings, so an unweighted cross-entropy would be
dominated by background the model classifies correctly from the first epoch,
while a fixed class weight would have to be retuned per recording. The focal
term instead adapts pixel by pixel, redirecting capacity to those that
remain ambiguous --- here the branch boundaries and the lowest-contrast
neurites. The
Dice term \cite{milletari2016vnet} compares the two masks as regions and
normalises by their combined area, which makes it independent of how much of
the frame the neuron occupies and supplies a gradient that depends on the
mask as a whole rather than on pixels independently --- both properties that
the per-pixel term lacks. Table~\ref{tab:loss} summarises the four.

\paragraph{Centreline-boundary Dice.}
cbDice \cite{shi2024cbdice} extends clDice \cite{shit2021cldice}, which
scores the skeletons of prediction and target against each other rather than
their masks: a skeleton holds far fewer pixels than the region it
summarises, so a break costing a handful of pixels of area removes a visible
fraction of the skeleton and the loss registers it. What clDice omits is
scale --- a centreline pixel inside the cell body and one on a two-pixel
branch count the same --- and cbDice supplies it by weighting each by its
distance to the nearest boundary.

Both quantities must be differentiable. The skeleton is obtained by
iterative morphological opening, following the clDice implementation:
erosion is a min-pool over the four-connected cross and dilation a max-pool
over the $3\times3$ square. Writing $I_k$ for the mask after $k$ erosions
and $S_0 = 0$,
\begin{equation}
S_k = \min\!\Big(1,\; S_{k-1} + \mathrm{ReLU}\!\big(I_k - \mathrm{open}(I_k)\big)\Big),
\qquad
\mathrm{open}(x) = \mathrm{dilate}\big(\mathrm{erode}(x)\big),
\label{eq:skel}
\end{equation}
accumulated over $k = 1,\dots,5$, the clamp keeping $S$ a soft membership in
$[0,1]$. Opening deletes exactly the thin ridges, so each iteration retains
the pixels that belong to thin structure and $S$ approximates the medial
axis. The weighting follows from the Euclidean distance transform: with
$d(m)$ the distance of each foreground pixel of a mask $m$ to its nearest
boundary,
\begin{equation}
W(m) = \frac{1}{d(m) + 1},
\qquad\text{normalised so that}\quad \textstyle\sum W = |m| ,
\label{eq:edtweight}
\end{equation}
which is $0.5$ on a one-pixel-radius branch and below $0.1$ inside the cell
body. The weight maps enter as constants: the distance transform is not
differentiated through, and $W_p$ is computed from the prediction
thresholded at $0.5$. For a predicted probability map $p$, a binary target
$g$, and their soft skeletons $S_p$ and $S_g$, weighted topological
precision and sensitivity are
\begin{equation}
\mathrm{cbT}_{\text{prec}}
 = \frac{\langle W_g \odot S_p,\; g \rangle}{\langle W_g,\; S_p \rangle},
\qquad
\mathrm{cbT}_{\text{sens}}
 = \frac{\langle W_p \odot S_g,\; p \rangle}{\langle W_p,\; S_g \rangle},
\label{eq:cbdice}
\end{equation}
and $\mathcal{L}_{\text{cbDice}} = 1 - 2\,\mathrm{cbT}_{\text{prec}}
\mathrm{cbT}_{\text{sens}} / (\mathrm{cbT}_{\text{prec}} +
\mathrm{cbT}_{\text{sens}})$. The two are asymmetric by construction:
precision asks whether the predicted centreline stays inside the target,
which detects spurious branches, and sensitivity whether the target
centreline stays inside the prediction, which detects breaks. A severed
branch is therefore penalised twice over --- the target centreline leaves
the prediction, and the weight at that location is at its largest --- and
the harmonic mean prevents either from being traded against the other. The
guarantee is partial: topology is preserved up to homotopy equivalence only
where the number of iterations exceeds the local radius
\cite{grim2025supervoxel}, so at $k = 5$ the term is exact along the
branches and approximate within the cell body, which is the regime we want.

\paragraph{Supervoxel connectivity.}
cbDice is a single soft score for the frame: it makes a break expensive
without identifying the pixels responsible, which is what the supervoxel
loss \cite{grim2025supervoxel} adds by testing each error component on its
own. Let $\hat{y}_-$ and $\hat{y}_+$ be the false-negative and
false-positive masks obtained by comparing the prediction $\hat{y}$ with the
target $y$. A connected component of $\hat{y}_-$ is \emph{negatively
critical} if it adjoins two or more distinct predicted components, so that
adding it would join them: it is a bridge, and without it the branch is
severed. A component of $\hat{y}_+$ is \emph{positively critical} if
deleting it separates predicted pixels that it had joined, which is decided
by removing it and relabelling; it falsely merges two structures. Every
other error component is a thickness error and is scored as one. Components
are four-connected throughout, and the prediction is binarised at $0.5$
before they are formed. Writing $\mathcal{N}$ and
$\mathcal{P}$ for the unions of the negatively and positively critical
components, $\Omega$ for the frame, and $\bar{\mathcal{L}}_0(A)$ for the
mean binary cross-entropy over a pixel set $A$,
\begin{equation}
\mathcal{L}_{\text{sv}}
 = (1-\alpha)\,\bar{\mathcal{L}}_0\big(\Omega \setminus (\mathcal{N} \cup \mathcal{P})\big)
 + \alpha\beta\,\bar{\mathcal{L}}_0(\mathcal{N})
 + \alpha(1-\beta)\,\bar{\mathcal{L}}_0(\mathcal{P}),
\label{eq:sv}
\end{equation}
with $\alpha = \beta = 0.5$, so that critical and non-critical pixels carry
equal total weight and splits and merges are penalised symmetrically. The
extra gradient thus lands on precisely the pixels whose presence or absence
flips a split into a join. Criticality is decided by connected-component
labelling, which is linear in the number of pixels and so far cheaper than
the persistent-homology formulations of the same idea. Both topology terms
are evaluated at $256\times256$ rather than at the full $1024\times1024$
output, since connectivity is a coarse property while the distance transform
and the component labelling dominate their cost; the critical masks are
upsampled before the cross-entropy is applied, so gradients are taken at
full resolution.

\paragraph{Optimisation.}
Training uses AdamW \cite{loshchilov2019adamw} with decoupled weight decay
$0.05$ and gradient-norm clipping at $1.0$, in mixed precision, for 110
epochs over clips of eight consecutive frames whose start points advance by
four, two clips per batch. All rates are warmed up over two epochs, since the
zero-initialised adapter matrices make the first steps sharp, then decayed
on a cosine schedule to $5\,\%$ of peak rather than to zero, which keeps the
late epochs productive on a dataset of this size. Augmentation is applied
identically to every frame of a clip --- horizontal and vertical flips
($p = 0.5$ each), rotation $\pm25^{\circ}$, shear $\pm20^{\circ}$,
brightness and contrast $\pm15\,\%$ --- so that the memory pathway is
trained against the motion actually present in the recording rather than
against apparent motion introduced by the augmentation. Validation runs
every second epoch and the best validation checkpoint is retained.

\begin{table}[!ht]
\centering
\caption{\textbf{The four terms of the training loss.} Each addresses a
distinct failure mode. The two topology terms are included because an error
that splits or merges a branch alters the area of the mask very little while
invalidating the morphology it is meant to describe.}
\label{tab:loss}
\small
\begin{tabular}{llcl}
\toprule
Term & Operates on & Weight & What it enforces \\
\midrule
Focal \cite{lin2017focal}        & per pixel            & 5.0 & robustness to severe, recording-dependent imbalance \\
Dice \cite{milletari2016vnet}    & region               & 2.0 & scale-invariant overlap on very small masks \\
cbDice \cite{shi2024cbdice}      & skeleton $+$ EDT     & 5.0 & thin-branch topology; reweights by inverse radius \\
Supervoxel \cite{grim2025supervoxel} & connected components & 0.3 & penalises splitting or merging a branch \\
\bottomrule
\end{tabular}
\end{table}
